Language models are defined as probability distributions over sequences of words, referred to as utterances. This chapter delves into the formalization of the term "utterance" and introduces fundamental concepts such as the alphabet, string, and language. Utilizing these concepts, a formal definition of a language model is presented, along with a discussion on the intricacies of defining distributions over infinite sets.

\subsection{Sets of Strings}
We begin by defining the very basic notions of alphabets and strings, where we take inspiration from \defn{formal language theory}.
First and foremost, formal language theory concerns itself with \emph{sets of structures}.
The simplest structure it considers is a \defn{string}.
So what is a string?
We start with the notion of an \defn{alphabet}.
\begin{definition}[Alphabet]
An \defn{alphabet} is a finite, non-empty set.
In this course, we will denote an alphabet using Greek capital letters, e.g., $\alphabet$ and $\outalphabet$.
We refer to the elements of an alphabet as \defn{symbols} or letters and will denote them with lowercase letters: $\syma$, $\symb$, $\symc$. 
\end{definition}

\begin{definition}[String]
A \defn{string}\footnote{A string is also referred to as a \defn{word}, which continues with the linguistic terminology.} over an alphabet is any \emph{finite} sequence of letters. 
Strings made up of symbols from $\alphabet$ will denoted by bolded Latin letters, e.g., $\stry = \sym_1 \cdots \sym_\strlen$ where each $\sym_n \in \alphabet$.
\end{definition}

The length of a string, written as $|\str|$, is the number of letters it contains. 
Usually, we will use $\strlen$ to denote $|\str|$ more concisely whenever the usage is clear from the context.
There is only one string of length zero, which we denote with the distinguished symbol $\eps$ and refer to as the \emph{empty string}.
By convention, $\eps$ is \emph{not} an element of the original alphabet.

New strings are formed from other strings and symbols with \defn{concatenation}.
Concatenation, denoted with $\strx \circ \stry$ or just $\strx\stry$, is an associative operation on strings.
Formally, the concatenation of two words $\str$ and $\strx$ is the word $\str \circ \strx = \str \strx$, which is obtained by writing the second argument after the first one. 
The result of concatenating with $\eps$ from either side results in the original string, which means that $\eps$ is the \defn{unit} of concatenation and the set of all words over an alphabet with the operation of concatenation forms a \defn{monoid}.
 
We have so far only defines strings as individual sequences of symbols.
To give our strings made up of symbols in $\alphabet$ a set to live in, we now define Kleene closure of an alphabet $\alphabet$.
\begin{definition}[Kleene Star]
Let $\alphabet$ be an alphabet.
The \defn{Kleene star} $\kleene{\alphabet}$ is defined as $\bigcup_{n=0}^\infty \alphabet^n$ where $\alphabet^n \defeq \underbrace{\alphabet \times \cdots \times \alphabet}_{n \text{ times}}$. 
Note that we define $\alphabet^0 \defeq \{\eps\}$.
We call the $\kleene{\alphabet}$ the \defn{Kleene closure} of the alphabet $\alphabet$.
\end{definition}
More informally, we can think of $\kleene{\alphabet}$ as the set which contains $\eps$ and all finite-length strings which can be constructed by concatenating arbitrary symbols from $\alphabet$. 
If the alphabet $\alphabet$ is empty, then its Kleene closure is $\left\{ \eps \right\}$. 
Otherwise, its Kleene closure is a \emph{countably infinite} set (this will come into play later!). 
It also leads us very naturally to our next definition.
\begin{definition}
Let $\alphabet$ be an alphabet. 
A \defn{language} $\lang$ is a subset of $\kleene{\alphabet}$.
\end{definition}
That is, a language is just a specified subset of all possible strings made up of the symbols in the alphabet.
This subset can be specified by simply enumerating a finite set of strings, or by a \emph{formal model}.
We will see examples of those later.
Importantly, these strings are \emph{finite}.
If not specified explicitly, we will often assume that $\lang = \kleene{\alphabet}$.

\paragraph{A note on terminology.} 
As we mentioned, these definitions are inspired by formal language theory.
We defined strings as our main structures of interest and symbols as their building blocks.
When we talk about natural language, the terminology is often slightly different: we may refer to the basic building blocks (symbols) as \defn{tokens} or \defn{words} (which might be composed of one or more \emph{characters} and form some form of ``words''; more on this in the chapter on tokenization) and their compositions (strings) as \defn{sequences} or \defn{sentences}.
Furthermore, what we refer to here as an alphabet may be called a \defn{vocabulary} (of words or tokens) in the context of natural language.
Sentences are therefore concatenations of words from a vocabulary in the same way that strings are concatenations of symbols from an alphabet.

\begin{example}[Kleene Closure]
Let $\alphabet = \left\{ \syma, \symb, \symc \right\}$. Then $$\kleene{\alphabet} = \{ \eps, \syma, \symb, \symc, \syma\syma, \syma \symb, \syma \symc, \symb\syma , \symb\symb, \symb\symc, \symc\syma , \symc\symb, \symc\symc, \syma \syma \syma , \syma \syma \symb, \syma \syma \symc, \dots \}.$$
Examples of a languages over this alphabet include $\lang_1 \defeq \set{\syma, \symb, \syma\symb, \symb\syma}$, $\lang_2 \defeq \set{\str \in \kleene{\alphabet} \mid \sym_1 = \syma}$, and $\lang_3 \defeq \set{\str \in \kleene{\alphabet} \mid |\str| \text{ is even}}$.
\end{example}

Next, we introduce two notions of subelements of strings. 
\begin{definition}[String Subelements]
A \defn{subsequence} of a string $\str$ is defined as a sequence that can be formed from $\str$ by deleting some or no symbols, leaving the order untouched. 
A \defn{substring} is a contiguous subsequence. 
For instance, $\syma\symb$ and $\symb\symc$ are substrings and subsequences of $\str=\syma\symb\symc$, while $\syma\symc$ is a subsequence but not a substring.
\defn{Prefixes} and \defn{suffixes} are special cases of substrings. 
A prefix is a substring of $\str$ that shares the same first letter as $\str$ and a suffix is a substring of $\str$ that shares the same last letter as $\str$.    
We will also denote a prefix $\sym_1\ldots\sym_{\idx - 1}$ of the string $\str = \sym_1\ldots\sym_\strlen$ as $\str_{< \idx}$.
\end{definition}

Finally, we also define the set of all \defn{infinite sequences} of symbols from some alphabet $\alphabet$ as $\alphabet^\infty$.
Since strings are canonically \emph{finite} in computer science, we will explicitly use the terms infinite sequence or infinite string to refer to elements of $\alphabet^\infty$.
Importantly, the set $\alphabet^\infty$ is \emph{uncountably infinite} for any $\alphabet$ such that $|\alphabet| \geq 2$, unlike $\kleene{\alphabet}$, which is countably infinite.

\subsection{Defining a Language Model}

We are now ready to introduce the main interest of the entire lecture series: language models.

\begin{definition}[Language model] \label{def:language-model}
Let $\alphabet$ be an alphabet.
A \defn{language model} is a (discrete) distribution $\pLM$ over $\kleene{\alphabet}$.
\end{definition}

\begin{example}[A very simple language model]
Let $\alphabet \defeq \set{\syma}$ and $\lang \defeq \kleene{\alphabet}$.
Define
\begin{equation*}
    \pLM\left(\syma^n\right) \defeq 2^{-\left(n + 1\right)},
\end{equation*}
where $\syma^0 \defeq \eps$ and $\syma^n \defeq \underbrace{\syma\ldots\syma}_{n \text{ times}}$.

We claim that $\pLM$ is a language model.
To see that, we verify that it is a valid probability distribution over $\lang$. 
It is easy to see that $\pLM\left(\syma^\idx\right) \geq 0$ for any $\idx$.
Additionally, we see that the probabilities of finite sequences indeed sum to $1$:
\begin{equation*}
    \sum_{\str \in \lang} \pLM\left(\str\right) = \sum_{\idx = 0}^{\infty} \pLM\left(\syma^\idx\right) 
    = \sum_{\idx = 0}^\infty 2^{-\left(\idx + 1\right)} 
    = \frac{1}{2}\sum_{\idx = 0}^\infty 2^{-\idx} 
    = \frac{1}{2} \frac{1}{1 - \frac{1}{2}} = 1.
\end{equation*}
\end{example}

A language model is itself a very simple concept---it is simply a distribution that weights strings (natural utterances) by their probabilities to occur in a particular language.
Note that we have not said anything about how we can represent or model this distribution yet.
Besides, for any (natural) language, the ground-truth language model $\pLM$ is of course \emph{unknown} and complex.
The next chapter, therefore, discusses in depth the computational models which we can use to try to tractably represent distributions over strings and ways of \emph{approximating} (learning) the ground-truth distribution based on finite datasets using such models. 
% \anej{Where do we put optimization, MLE, backprop...?\response{clara} thats in the parameter estimation section \response{anej} Yup, thanks so much! :D}
